% =============================================================
% Tools of the Trade
% Chapter order is set in ../main.tex; the rendered chapter number
% comes from \chapter{} below.
%
% Stub chapter. Includes the prompting-craft material that was
% previously sketched in manuscript/prompting.md, plus other
% practical patterns (output parsing, retry / fallback, common
% chains and idioms). Fleshing out is deferred.
% =============================================================

\chapter{Tools of the Trade}
\label{ch:tools}

\todo{Stub chapter. Fleshing out is deferred. The outline below
names the sections this chapter is intended to cover; each will
become several pages of practical material once the foundational
chapters (\ref{ch:mental} and \ref{ch:working}) are complete.}

This chapter is the practical-skills counterpart to the operational
material of Chapter~\ref{ch:working}. Where Chapter~\ref{ch:working}
covered \emph{how to call a model}, this one covers \emph{how to
get useful work out of one}. The boundary is the system prompt: on
one side, the API contract that tells the model what tools it has
and what role it is playing; on the other, the craft of phrasing
that makes the difference between a model that writes the code you
asked for and a model that writes plausible-looking code that does
not compile. Both sides of the boundary matter; we take the second
one here.

\section{Prompting Patterns}
\label{sec:tools-prompting}

\todo{The substantive section. Cover, with worked examples in the
Gemma~4 / Ollama setup used in Chapter~\ref{ch:toolcalls}: the
core prompting moves (zero-shot, few-shot, chain-of-thought, role
prompting); when each helps and when each fails; the prompting
oddities the literature has documented (capitalization sensitivity,
the prompt-repetition effect, sensitivity to the order of few-shot
examples). Quote each piece of advice from a primary source so the
material does not turn into received wisdom.}

\section{Context Construction}
\label{sec:tools-context}

\todo{The observation that constructing the right \emph{context}
matters more than crafting the right \emph{prompt}. What goes in
the system prompt vs the user message vs the tool result. Treat
prompting as one lever among several --- alongside what
information to include, in what order, with what compression --- in
shaping the input the model sees.}

\section{Output Parsing and Recovery}
\label{sec:tools-parsing}

\todo{The model produced something; now what? Defensive parsing
patterns; common failure modes (markdown fences around JSON,
trailing prose, hallucinated keys); soft-recovery loops; when to
fail-fast and when to retry. Connection to structured outputs
(\S\ref{sec:structured}).}

\section{Retry and Fallback}
\label{sec:tools-retry}

\todo{The agent loop will fail; the question is what happens
next. Retry strategies (with and without back-off); model
fallback (start with a small model, escalate to a larger one on
failure); deterministic fallback to a non-LLM path for known-bad
inputs.}

\section{Common Idioms}
\label{sec:tools-idioms}

\todo{A short menagerie of patterns that recur: the
classify-then-route pattern, the propose-and-verify pattern, the
reduce-with-summary pattern, the self-critique pattern. Each gets
a paragraph naming the shape and the case it solves.}
